\section{Mandat d'investissement et philosophie}
\label{sec:mandate}

\subsection{Mandat}

\begin{center}
\renewcommand{\arraystretch}{1.35}
\begin{tabular}{p{4.2cm}p{10cm}}
\toprule
\textbf{Univers d'investissement} & Paires de devises majeures~: EUR-USD, GBP-USD, USD-JPY, USD-CAD, plus les deux métaux précieux~: l'or (XAU-USD) et l'argent (XAG-USD). Pas d'exposition aux devises émergentes, aux \emph{crosses} exotiques ni aux dérivés structurés. \\
\textbf{Objectif de rendement} & CAGR cible de l'ordre de \textbf{40\,\%} annualisés. \\
\textbf{Contrainte de risque} & \textbf{Aucune}. Le plafond de drawdown de 35\,\% des versions antérieures a été retiré sur instruction du mandant, qui a déclaré le risque non contraignant. Le repli est mesuré et rapporté, il ne borne pas le dimensionnement. \\
\textbf{Horizon} & Multi-annuel (minimum 3~ans recommandé pour réalisation statistique des métriques). \\
\textbf{Cadre réglementaire} & L'exécution suppose un compte autorisant le levier requis par le dimensionnement retenu. La compatibilité avec les plafonds ESMA retail (30$\times$ sur majors G10) doit être vérifiée auprès du courtier avant tout déploiement. \\
\textbf{Réserve de capital} & À dimensionner sur le repli d'équité constaté (50,9\,\%), et non sur un plafond contractuel qui n'existe plus. \\
\textbf{Rebalancement} & Quotidien, sans adaptation de régime. Pas de réallocation discrétionnaire. \\
\textbf{Liquidité} & Positions entièrement liquidables sous 24~heures. \\
\textbf{Statut actuel} & Prêt pour déploiement contrôlé en \emph{paper-trade}. Go-live conditionné à la validation de la check-list de production (section~\ref{sec:production}). \\
\bottomrule
\end{tabular}
\end{center}

\begin{insightbox}
\textbf{Ce mandat a changé de nature.} Les versions antérieures de ce document visaient 10 à 15\,\% de rendement sous un plafond de repli de 35\,\%~: l'objectif était la régularité. Le mandat actuel vise 40\,\% sans contrainte de repli. Ce n'est pas le même produit, et les deux ne se comparent pas. Tout ce qui suit décrit le second.
\end{insightbox}

\subsection{Problème adressé}

Le marché des changes est simultanément l'un des plus liquides et l'un des plus difficiles à battre systématiquement. Chaque signal pris isolément est dominé par un bruit important, et les \emph{edges} individuels tendent à s'éroder rapidement lorsqu'ils sont répliqués par d'autres acteurs. La réponse classique à ce problème~--- empiler plusieurs signaux indépendants pour réduire la variance du portefeuille agrégé~--- n'est efficace qu'à une condition stricte~: les signaux doivent être \textbf{vraiment orthogonaux} les uns aux autres. Or, la plupart des stratégies FX publiées sont en réalité fortement corrélées parce qu'elles exploitent la même source d'alpha sous-jacente (\emph{carry}, tendance EMA, \emph{mean reversion} RSI).

\begin{thesisbox}
L'ambition de la Stratégie FX Multi-Moteurs est d'assembler des moteurs de rendement qui capturent des \textbf{sources d'alpha structurellement différentes}, de sorte que leur combinaison porte le rendement visé sans concentrer le risque sur un seul régime de marché. La section~\ref{sec:limitations} montre dans quelle mesure cette ambition est tenue --- et dans quelle mesure elle ne l'est pas.
\end{thesisbox}

\subsection{Architecture du portefeuille}

\begin{figure}[htbp]
\centering
\small
\begin{tabular}{cccc}
\colorbox{softBlue}{\parbox{3.3cm}{\centering\color{primary}\sffamily\textbf{Moteur 1\\Mean Reversion Intraday}\\[0.2em]\footnotesize 62\,\% de l'allocation\\Retour au VWAP filtré macro\\4~paires, minute}} &
\colorbox{softYellow}{\parbox{3.3cm}{\centering\color{accent}\sffamily\textbf{Moteur 2\\Momentum Journalier}\\[0.2em]\footnotesize 9\,\% de l'allocation\\EMA 20/50 + confirmation RSI\\3~paires}} &
\colorbox{softGreen}{\parbox{3.3cm}{\centering\color{rsiGreen}\sffamily\textbf{Moteur 3\\RSI Journalier}\\[0.2em]\footnotesize 9\,\% de l'allocation\\Mean reversion sur RSI(14)\\3~paires}} &
\colorbox{lightGrey}{\parbox{3.3cm}{\centering\color{apogeeGold}\sffamily\textbf{Moteur 4\\Momentum Multi-Actifs}\\[0.2em]\footnotesize 20\,\% de l'allocation\\Suivi de tendance\\XAU-USD, USD-JPY, XAG-USD}} \\[1.2em]
\multicolumn{4}{c}{$\Downarrow$} \\[0.3em]
\multicolumn{4}{c}{\colorbox{lightGrey}{\parbox{12cm}{\centering\color{primary}\sffamily\textbf{Allocation statique 62 / 9 / 9 / 20}\\[0.1em]\footnotesize Pondérations fixes, pas de réallocation discrétionnaire}}} \\[0.6em]
\multicolumn{4}{c}{$\Downarrow$} \\[0.3em]
\multicolumn{4}{c}{\colorbox{lightGrey}{\parbox{12cm}{\centering\color{primary}\sffamily\textbf{Couche de ciblage de volatilité}\\[0.1em]\footnotesize Vol cible 0.37, levier maximal 31$\times$, estimateur pessimiste anti look-ahead\\Plafond de drawdown désactivé}}} \\[0.6em]
\multicolumn{4}{c}{$\Downarrow$} \\[0.3em]
\multicolumn{4}{c}{\colorbox{primary}{\parbox{12cm}{\centering\color{white}\sffamily\textbf{Équité finale du portefeuille}}}} \\
\end{tabular}
\caption{Architecture à quatre moteurs de la Stratégie FX Multi-Moteurs. Allocation statique, puis couche de ciblage de volatilité, puis pipeline d'équité finale. Les 10~points de capital transférés à la sleeve momentum en juillet~2026 ont été pris sur la seule Mean Reversion Intraday~: les poids absolus du Momentum Journalier et du RSI Journalier sont inchangés, si bien que le trio FX répartit entre ses membres les 80\,\% qui lui restent selon les ratios $77/11/11$.}
\end{figure}

\subsection{Pourquoi quatre moteurs}

\begin{itemize}
\item \textbf{Deux moteurs ne suffisent pas}~: les analyses d'attribution montrent qu'une combinaison à deux moteurs (Mean Reversion Intraday + Momentum Journalier) laisse les années 2019 et 2023 en territoire négatif.
\item \textbf{Le trio FX seul plafonne}~: sa combinaison non leviérée rend 2,2\,\% pour une volatilité de 5,8\,\%. Poussé par le ciblage de volatilité, son rendement culmine autour de 32\,\% avant que le \emph{drag} de variance ne le fasse retomber. Le mandat de 40\,\% lui était donc inaccessible, quel que soit le levier.
\item \textbf{La sleeve momentum lève cette contrainte parce qu'elle est quasi orthogonale au trio}~: à 20\,\% de poids, elle porte le Sharpe du combiné de 0,66 à environ 1,0 (tableau~\ref{tab:robustness_bootstrap_ci}), ce qui déplace le plafond de rendement au-delà de la cible. La cible devient atteignable à un niveau de volatilité \emph{inférieur} à celui qu'exigeait le trio seul.
\item \textbf{Trois instruments plutôt qu'un depuis juillet~2026}~: la sleeve applique le même signal de tendance à l'or, au dollar-yen et à l'argent, chacun recevant un tiers du budget. L'argent, médiocre pris isolément, casse la dépendance du portefeuille à la seule trajectoire de l'or~; le dollar-yen y ajoute un \emph{carry} de portage positif, alors que l'or coûte des frais de portage.
\item \textbf{Un cinquième signal a été évalué et retiré}~: une composante \emph{cross-section momentum} présentait une corrélation supérieure à 0,3 avec le Momentum Journalier, sans apport au Sharpe combiné.
\end{itemize}

\begin{insightbox}
\textbf{La leçon principale de la recherche préparatoire~:} un moteur avec un Sharpe standalone modeste peut être le meilleur diversificateur si sa corrélation avec le reste du portefeuille est négative pendant les années difficiles. Ce qui compte n'est pas le Sharpe individuel mais la \textbf{contribution marginale au Sharpe du portefeuille}, qui dépend de la corrélation. Un signal à Sharpe 0.16 anti-corrélé bat souvent un signal à Sharpe 0.60 corrélé.
\end{insightbox}

\begin{insightbox}
\textbf{Le revers de cette leçon apparaît à l'exécution.} L'orthogonalité qui justifie l'entrée de la sleeve momentum explique aussi qu'elle porte, à elle seule, 91\,\% du résultat du portefeuille sur la période mesurée~: la diversification a servi à \emph{autoriser} le rendement visé, pas à le répartir. Le passage de l'or seul au trio d'instruments a diversifié la sleeve, pas le portefeuille~--- la concentration s'est même accentuée avec son poids. Voir la section~\ref{sec:limitations}.
\end{insightbox}
